\section{Introduction}
\label{sec:intro}

Mobile apps have become the primary surface through which people manage their daily lives, from sending messages to tracking expenses to navigating unfamiliar cities. A mobile GUI agent is a system that operates these apps on a user's behalf: it reads what is on the screen, decides on a tap, swipe, or text input, and observes the result, repeating until the user's request is fulfilled. With strong vision-language models behind them, recent agents have begun to handle multi-step requests reliably on a fixed set of apps, opening the door to personal assistants that act, not just chat.

The past two years have produced a wave of capable GUI agents~\citep{cheng2024seeclick,hong2024cogagent,yang2023appagent,guo2025agentsama}. SeeClick~\citep{cheng2024seeclick} introduced grounding-aware pretraining for screen elements, CogAgent~\citep{hong2024cogagent} scaled a GUI-specialized vision-language model with high-resolution support, and AppAgent~\citep{yang2023appagent} recorded UI documentation during an offline exploration phase to guide deployment. Benchmarks have advanced in parallel. AndroidWorld~\citep{rawles2024androidworld} provides programmatic tasks with reward signals derived from the device state. AndroidLab~\citep{xu2024androidlab} adds a reproducible training framework, AndroidControl~\citep{li2024androidcontrol} contributes a 15K-demonstration dataset over 833 apps, and MobileWorld~\citep{mobileworld2025} extends task horizons and introduces multi-application and MCP-augmented categories. The best systems now reach 60--80\% success rates on these benchmarks~\citep{openmobile2026,mobilerl2025}.

However, headline numbers like these typically come from evaluations where training and test tasks share the same template family, and held-out conditions are defined at the parameter or seed level rather than at the app or category level. A 2,000-task training set built by varying the parameters of 116 evaluation templates is, in agent terms, the same task with different inputs---not a measure of how the agent handles something genuinely new. The deployment-relevant question is different. When a user installs a regional banking app, a niche notes tool, or a ride-hailing service the agent has never encountered, can the agent become reliable on that app from a handful of adaptation examples, possibly without any analog in its pretraining? This is the gap we focus on, and current evaluation practice does not measure it directly.

The cost structure of mobile agent deployment makes this question pressing. Synthesizing thousands of task variants and fine-tuning per app is not viable when a user installs three new apps in a week. Online RL with large training budgets is not viable when a user expects the agent to start working within minutes. The natural compute budget at deployment time is on the order of a few dozen trajectories on a small adaptation set, and an effective agent has to make that budget count. Two ideas point at how. The first is to give the agent a piece of structured prior knowledge that transfers across apps, so adaptation starts from something better than scratch. The second is to update the agent's weights during adaptation, using whatever reward signal it can extract from its own rollouts on the small adaptation set. Recent work in single-turn reasoning has shown that doing both jointly outperforms either alone~\citep{zhang2026espl}, but whether this carries over to the multi-step, visual, expensive-rollout GUI agent setting under strict cross-app evaluation is unanswered.

We introduce \textbf{EvoFSM-RL}, a framework that combines both ideas in a way designed for the deployment-time budget. The agent's prior knowledge is a two-layer finite state machine: the lower layer captures app-specific UI structure, and the upper layer captures category-generic workflow patterns written without any app-specific tokens. A linter enforces this separation, so the upper layer is the structurally guaranteed unit of cross-app transfer. Same-category source apps aggregate their upper layers into a per-category abstract library $L_C$, retrieved by Play Store category at deployment time. Adaptation runs a joint loop: a frozen reference language model proposes structured edits to the upper-layer FSM, and a LoRA adapter on the policy is updated by group-relative policy optimization with within-(FSM, task) advantage baselines. The same loop runs at source-pool pretraining (producing a shared LoRA $\pi^{\text{pre}}_\theta$) and at target-app deployment (continuing from $\pi^{\text{pre}}_\theta$), making ``test-time adaptation'' a well-defined operation rather than an ad-hoc fine-tune.

To measure cross-app and cross-category generalization honestly, we evaluate on a 25-app Android benchmark with three independent splits stacked on top of one another: source pool vs.\ target apps (pool-disjoint), near-transfer vs.\ far-transfer target categories (category-disjoint), and within-app $T_{\text{adapt}}$ vs.\ $T_{\text{eval}}$ templates (template-disjoint). Against a four-rung ablation, static category knowledge alone (B2) lifts success rate by $+9.3$~pp over a zero-shot baseline (B1) on near-transfer apps, online upper-layer evolution (B3) adds another $+3.7$~pp on top, and the far-transfer panel serves as a built-in null control on the injection mechanism. The full joint variant (B4) is the test of whether procedural and declarative knowledge can be co-adapted at deployment-time scale.

\medskip
\noindent\textbf{Contributions.} This paper makes the following contributions:
\begin{itemize}
    \item \textbf{A two-layer FSM as structured agent knowledge.} We propose a representation in which an app-specific lower layer and a category-generic upper layer are separated by an enforced lint pass. The upper layer is the structurally guaranteed unit of cross-app transfer and the only target of online mutation.

    \item \textbf{The EvoFSM-RL adaptation loop.} A joint evolution-plus-RL procedure that mutates the upper-layer FSM via a frozen reference language model and updates a LoRA adapter via group-relative policy optimization on within-(FSM, task) baselines. The same loop runs at source-pool pretraining and at target-app deployment, with $L_C$ mediating between them.

    \item \textbf{An evaluation protocol disjoint on three independent levels.} Pool-disjoint between source and target apps, category-disjoint between near and far transfer, and template-disjoint between adaptation and evaluation within each target app, giving an honest measure of generalization rather than within-template parameter robustness.

    \item \textbf{Empirical evidence and failure-mode characterization.} The four-rung ablation isolates the contributions of static category knowledge, online upper-layer evolution, and joint weight adaptation, with consistent gains on near-transfer apps and a null control on far-transfer. We also characterize when the joint mechanism succeeds (base policy with nonzero success rate on $T_{\text{adapt}}$) and when it does not (uniformly failing rollouts collapse the GRPO advantage signal), pointing to concrete follow-up directions.
\end{itemize}
